\documentclass[11pt,a4paper]{article}

\usepackage[utf8]{inputenc}
\usepackage[T1]{fontenc}
\usepackage[spanish,es-noquoting]{babel}
\usepackage{amsmath,amssymb,amsthm}
\usepackage{graphicx}
\usepackage{booktabs}
\usepackage{microtype}
\usepackage[margin=2.7cm]{geometry}
\usepackage{xcolor}
\usepackage[colorlinks=true,linkcolor=azulosc,citecolor=azulosc,urlcolor=azulosc]{hyperref}

\definecolor{azulosc}{HTML}{1F4E79}
\definecolor{grisclaro}{HTML}{F2F2F2}

\usepackage{tcolorbox}
\tcbuselibrary{skins,breakable}

\newtcolorbox{idea}[1][]{
  colback=grisclaro, colframe=azulosc, boxrule=0.6pt, arc=2pt,
  left=6pt, right=6pt, top=5pt, bottom=5pt, breakable, #1}

\newtcolorbox{aviso}[1][]{
  colback=white, colframe=black!45, boxrule=0.6pt, arc=2pt,
  left=6pt, right=6pt, top=5pt, bottom=5pt, breakable, #1}

\newcommand{\dd}{\mathrm{d}}
\newcommand{\ee}{\mathrm{e}}
\newcommand{\ii}{\mathrm{i}}
\newcommand{\rp}{p_r}
\newcommand{\Jr}{J}
\newcommand{\Qr}{Q}

\theoremstyle{definition}
\newtheorem{ejercicio}{Ejercicio}

\title{\bfseries Phase mixing, amortiguamiento de Landau y variables\\
       ángulo--acción: una introducción computacional\\
       al sistema de Vlasov--Poisson en simetría esférica}
\author{Notas de trabajo}
\date{}
\begin{document}
\maketitle

\begin{abstract}
\noindent
Estas notas explican, desde cero y con detalle, cómo se simula numéricamente un
sistema estelar no colisional y cómo se distinguen dos mecanismos que hacen
desaparecer una perturbación sin que haya disipación alguna: el \emph{phase
mixing}, que es puramente cinemático, y el \emph{amortiguamiento de Landau}, que
es colectivo. Se presenta el sistema de Vlasov--Poisson en simetría esférica con
momento angular fijo, las variables ángulo--acción del potencial isócrono, el
método \emph{particle-in-cell} con integradores simplécticos, y un diagnóstico
espectral $h_k(t)$ que permite medir el mezclado. Todo se ilustra con resultados
de un código real, incluidos los errores cometidos durante su desarrollo, porque
la forma en que se descartan las explicaciones falsas es parte del método. El
hilo conductor es un caso de estudio: una meseta inesperada en $h_1(t)$ al
encender la autogravedad, que tras cuatro barridos de control resulta no ser ni
ruido numérico ni un modo colectivo, sino la huella de estar usando las
coordenadas equivocadas.
\end{abstract}

\tableofcontents
\bigskip

\section{Por qué este problema es interesante}

Un cúmulo estelar, una galaxia o un halo de materia oscura son sistemas donde las
colisiones entre partículas son tan raras que pueden ignorarse durante la vida del
universo. La dinámica está dominada por el campo gravitatorio \emph{colectivo} que
todas las partículas generan juntas. El objeto matemático que describe esto es la
ecuación de Vlasov acoplada a la de Poisson.

Lo que hace al sistema fascinante, y lo que estas notas persiguen, es que
\textbf{una perturbación puede desaparecer sin que se pierda información}. No hay
viscosidad, no hay colisiones, la entropía de grano fino se conserva exactamente.
Y sin embargo, si uno mide casi cualquier cantidad macroscópica, la ve decaer.
Hay dos mecanismos distintos detrás de esto y conviene no confundirlos:

\begin{idea}
\textbf{Phase mixing.} Órbitas vecinas tienen frecuencias ligeramente distintas.
Una estructura inicial se \emph{enrolla} en el espacio de fases hasta hacerse
más fina que cualquier resolución. La información sigue ahí, pero deja de ser
accesible a cualquier medida promediada. Es un efecto \textbf{cinemático}: ocurre
incluso con partículas de prueba en un potencial fijo.

\medskip
\textbf{Amortiguamiento de Landau.} La perturbación genera un campo, el campo
actúa sobre las partículas, y las que orbitan en resonancia con la perturbación
intercambian energía con ella de forma neta. Es un efecto \textbf{colectivo}:
desaparece si se apaga la autogravedad.
\end{idea}

Distinguirlos experimentalmente es el objetivo central de estas notas, y se puede
hacer de forma tajante: como se demuestra en la sección~\ref{sec:envolvente}, el
phase mixing produce una envolvente \emph{gaussiana} en el tiempo y el
amortiguamiento de Landau una \emph{exponencial}. Basta ajustar el logaritmo de la
amplitud contra $t^2$ y contra $t$ y ver cuál gana.

\subsection*{A quién van dirigidas}

Se supone familiaridad con mecánica hamiltoniana, análisis de Fourier y métodos
numéricos básicos. No se supone nada de dinámica estelar ni de métodos de
partículas. Los apartados marcados como \textsc{detalle numérico} pueden saltarse
en una primera lectura sin perder el hilo físico.

\section{El sistema}

\subsection{Vlasov--Poisson}

Sea $f(\mathbf{x},\mathbf{v},t)$ la densidad en el espacio de fases: $f\,\dd^3x\,
\dd^3v$ es la masa contenida en un elemento de volumen. Si no hay colisiones, $f$
se conserva a lo largo de las trayectorias, lo que es la ecuación de Vlasov (o de
Boltzmann sin colisiones):
\begin{equation}
  \frac{\partial f}{\partial t}
  + \mathbf{v}\cdot\frac{\partial f}{\partial \mathbf{x}}
  - \nabla\Phi\cdot\frac{\partial f}{\partial \mathbf{v}} = 0 ,
  \label{eq:vlasov}
\end{equation}
acoplada al potencial que la propia masa genera,
\begin{equation}
  \nabla^2\Phi = 4\pi G \rho, \qquad
  \rho(\mathbf{x},t) = \int f(\mathbf{x},\mathbf{v},t)\,\dd^3 v .
  \label{eq:poisson}
\end{equation}

La ecuación~\eqref{eq:vlasov} dice simplemente que $f$ es constante siguiendo el
flujo hamiltoniano. Es una ecuación de transporte, no difusiva: \emph{nada} en
ella destruye información. Toda la riqueza está en que el campo $\Phi$ depende de
$f$, lo que la vuelve no lineal e integro-diferencial.

\subsection{Reducción a un grado de libertad}

Trabajaremos en simetría esférica. Entonces el momento angular $L=|\mathbf{x}
\times\mathbf{v}|$ de cada partícula se conserva, y el movimiento radial se separa.
Tomamos además una distribución con un \emph{único} valor del momento angular,
\begin{equation}
  f(\mathbf{x},\mathbf{v},t) \;\propto\; F(r,\rp,t)\,\delta(L-L_0),
  \label{eq:deltaL}
\end{equation}
lo que reduce el problema a un grado de libertad con hamiltoniano
\begin{equation}
  H(r,\rp) = \frac{\rp^2}{2} + \Phi_{\rm ef}(r),
  \qquad
  \Phi_{\rm ef}(r) = \Phi(r) + \frac{L_0^2}{2r^2}.
  \label{eq:hamiltoniano}
\end{equation}

Esta elección no es una limitación caprichosa: un sistema con un grado de libertad
y potencial acotado es \emph{siempre integrable}, luego admite variables
ángulo--acción exactas. Eso nos da una solución de referencia contra la que medir,
que es exactamente lo que hace falta para validar un código.

\begin{aviso}
\textbf{Unidades.} Tomamos $G=1$ y la masa y el radio de escala del potencial de
fondo iguales a uno. Con ello el tiempo orbital característico resulta del orden
de $10^2$, y las masas de la componente que perturba se expresarán como fracción
de la del fondo. En todas las corridas de estas notas $L_0=2$.
\end{aviso}

\section{Variables ángulo--acción}
\label{sec:aa}

\subsection{Qué son y por qué se usan}

Para un sistema integrable existe una transformación canónica $(r,\rp)\to(\Qr,\Jr)$
tal que el hamiltoniano depende \emph{solo} de la acción:
\begin{equation}
  H = H(\Jr) \quad\Longrightarrow\quad
  \dot{\Jr} = -\frac{\partial H}{\partial \Qr} = 0,
  \qquad
  \dot{\Qr} = \frac{\partial H}{\partial \Jr} \equiv \omega(\Jr).
  \label{eq:ecmov}
\end{equation}
La acción se define como el área encerrada en el espacio de fases,
\begin{equation}
  \Jr = \frac{1}{2\pi}\oint \rp\,\dd r
      = \frac{1}{\pi}\int_{r_-}^{r_+}\sqrt{2\big(E-\Phi_{\rm ef}(r)\big)}\;\dd r,
  \label{eq:accion}
\end{equation}
donde $r_\pm$ son los puntos de retorno, raíces de $E=\Phi_{\rm ef}(r)$, y el
ángulo $\Qr\in[0,2\pi)$ avanza uniformemente con el tiempo.

La consecuencia es que \textbf{la dinámica se vuelve trivial}: la solución exacta
de~\eqref{eq:vlasov} sin autogravedad es
\begin{equation}
  \boxed{\;F(\Qr,\Jr,t) = F_0\big(\Qr - \omega(\Jr)\,t,\;\Jr\big)\;}
  \label{eq:solucion}
\end{equation}
Cada valor de la acción gira a su propia velocidad angular. Eso es todo el phase
mixing, escrito en una línea.

\subsection{El potencial isócrono}

El potencial isócrono, introducido por Hénon~\cite{henon1959}, es especial: es el
\emph{único} potencial esférico, además del armónico y el kepleriano, para el que
las acciones se obtienen en forma cerrada. En nuestras unidades,
\begin{equation}
  \Phi_{\rm iso}(r) = -\frac{1}{1+\sqrt{1+r^2}} .
  \label{eq:isocrono}
\end{equation}
Interpola entre un núcleo armónico ($\Phi\simeq -\tfrac12 + \tfrac14 r^2$ para
$r\ll1$) y un exterior kepleriano ($\Phi\simeq -1/r$ para $r\gg1$), lo que lo hace
un modelo de galaxia razonable y a la vez manejable.

Su acción radial es~\cite[\S3.5.2]{binney2008}
\begin{equation}
  \Jr(E,L) = \frac{1}{\sqrt{-2E}} - \frac{1}{2}\Big(L+\sqrt{L^2+4}\Big),
  \label{eq:Jiso}
\end{equation}
e invirtiendo, con $c\equiv\tfrac12\big(L+\sqrt{L^2+4}\big)$,
\begin{equation}
  E(\Jr) = -\frac{1}{2(\Jr+c)^2},
  \qquad
  \omega(\Jr) = \frac{\partial E}{\partial \Jr} = \frac{1}{(\Jr+c)^3},
  \qquad
  \omega'(\Jr) = -\frac{3}{(\Jr+c)^4}.
  \label{eq:omega}
\end{equation}

Las dos últimas expresiones son el corazón de todo lo que sigue. La
\textbf{dispersión de frecuencias} $\omega'$ es la que produce el enrollamiento:
si $\omega'$ fuese cero, todas las órbitas girarían al unísono y no habría phase
mixing en absoluto.

Para $L_0=2$ resulta $c=1+\sqrt{2}\simeq2.4142$, y sobre el soporte típico de
nuestras distribuciones ($\Jr\in[0,\,0.35]$) las frecuencias recorren la banda
\begin{equation}
  \omega \in [0.0473,\;0.0711],
  \qquad\text{periodo orbital}\;\; T=2\pi/\omega \simeq 100.
  \label{eq:banda}
\end{equation}
Esa banda reaparecerá en la sección~\ref{sec:autogravedad} como el criterio que
decide si un modo colectivo es resonante o no.

\subsection{El mapa explícito, y una ecuación tipo Kepler}
\label{sec:mapa}

Para pasar de $(r,\rp)$ a $(\Qr,\Jr)$ en el isócrono se procede así. Dada la
energía $E$ del hamiltoniano~\eqref{eq:hamiltoniano}, los puntos de retorno son
\begin{equation}
  r_\pm^2 = \frac{1+E(2+L^2)\pm\sqrt{1+2E\,(2+2E+L^2)}}{2E^2},
\end{equation}
la acción sale de~\eqref{eq:Jiso}, y el ángulo se obtiene a través de una variable
auxiliar $\eta$ (el análogo de la anomalía excéntrica de Kepler). Con
$s\equiv 1+\sqrt{1+r^2}$ y $s_\pm$ sus valores en los puntos de retorno,
\begin{equation}
  \cos\eta = \frac{s_-+s_+-2s}{s_+-s_-},
  \qquad
  \Qr = \eta - e\,\sin\eta,
  \label{eq:kepler}
\end{equation}
con una excentricidad efectiva $e$ que depende de $E$ y $L$. El signo de $\rp$
decide en qué media órbita estamos.

\begin{aviso}
\textsc{detalle numérico.} El mapa \emph{directo} $(r,\rp)\to(\Qr,\Jr)$ es
explícito: basta evaluar~\eqref{eq:kepler}. El \emph{inverso} no lo es, porque hay
que resolver $\Qr = \eta - e\sin\eta$ para $\eta$. Es la ecuación de Kepler, y se
resuelve con Newton--Raphson en unas cinco iteraciones. En el código esto aparece
al construir rejillas de cuadratura directamente en $(\Qr,\Jr)$.
\end{aviso}
\section{Phase mixing: la teoría}
\label{sec:mixing}

\subsection{El observable}

Necesitamos una cantidad escalar que mida cuánta estructura angular queda. La
elección natural es proyectar $F$ sobre una función de prueba $\Phi_{\rm p}$ y
quedarse con los armónicos de Fourier en el ángulo:
\begin{equation}
  h_k(t) \;=\; 8\pi^2 L_0 \iint F(\Qr,\Jr,t)\;\Phi_{\rm p}(\Qr,\Jr)\;
             \ee^{-\ii k \Qr}\,\dd \Qr\,\dd \Jr .
  \label{eq:hk}
\end{equation}
El prefactor $8\pi^2L_0$ es la medida de los grados de libertad angulares que la
delta de~\eqref{eq:deltaL} colapsa; con él, $h_0$ es literalmente la masa cuando
$\Phi_{\rm p}=1$.

Tomamos una función de prueba separable, $\Phi_{\rm p}=A(\Qr)\,B(\Jr)$, con
\begin{equation}
  A(\Qr)=\exp\!\Big[-\frac{\sin^2(\Qr/2)}{\sigma_Q^2}\Big],
  \qquad
  B(\Jr)=\exp\!\Big[-\frac{(\Jr-\Jr_0)^2}{\sigma_J^2}\Big]\,\Jr^2 .
  \label{eq:prueba}
\end{equation}
El perfil en $\Qr$ es periódico y par (una gaussiana «envuelta», emparentada con la
distribución de von Mises), de modo que sus coeficientes de Fourier $a_k$ son
reales.

\begin{idea}
\textbf{Por qué $h_k$ y no otra cosa.} Porque su evolución bajo phase
mixing se calcula \emph{exactamente}, y por tanto sirve de patrón contra el que
medir un código. Cualquier otra cantidad macroscópica se comportaría
cualitativamente igual, pero sin una respuesta cerrada con la que compararla.
\end{idea}

\subsection{La solución cerrada}

Sustituyendo~\eqref{eq:solucion} en~\eqref{eq:hk}, y escribiendo los coeficientes
de Fourier en $\Qr$ de la distribución inicial,
\begin{equation}
  c_k(\Jr) = \frac{1}{2\pi}\int_0^{2\pi} F_0(\Qr,\Jr)\,\ee^{-\ii k\Qr}\,\dd \Qr ,
\end{equation}
la integral angular se hace de una vez y queda
\begin{equation}
  \boxed{\;
  h_k(t) \;=\; a_0\,a_k\;
  \frac{\displaystyle\int \dd \Jr\; B(\Jr)\,c_k(\Jr)\,\ee^{-\ii k\,\omega(\Jr)\,t}}
       {\displaystyle\int \dd \Jr\; c_0(\Jr)} \;}
  \label{eq:hkexacto}
\end{equation}
donde $a_0$ es la masa total. Esta fórmula vale para \emph{cualquier} $F_0$,
separable o no; si no es separable, $c_k$ simplemente depende de $\Jr$ y es
compleja.

Nótese que $h_0$ no contiene el factor oscilante: \textbf{$h_0$ es constante en el
tiempo}. El phase mixing no destruye masa, solo la enrolla. Eso da una
comprobación gratuita en toda simulación, y en la sección~\ref{sec:autogravedad}
se convertirá en el diagnóstico más informativo de todos.

\subsection{Por qué la envolvente es gaussiana}
\label{sec:envolvente}

Supongamos que el integrando de~\eqref{eq:hkexacto} está concentrado alrededor de
$\Jr_0$ con anchura $\sigma$, es decir $B\,c_k \propto
\exp[-(\Jr-\Jr_0)^2/\sigma^2]$. Linealizamos la frecuencia,
\begin{equation}
  \omega(\Jr) \simeq \omega_0 + \omega'\cdot(\Jr-\Jr_0),
\end{equation}
y usamos la transformada de Fourier de una gaussiana,
\begin{equation}
  \int_{-\infty}^{\infty} \exp\!\big(-x^2/\sigma^2\big)\,\ee^{-\ii a x}\,\dd x
  = \sqrt{\pi}\,\sigma\,\exp\!\big(-a^2\sigma^2/4\big),
\end{equation}
con $a=k\,\omega'\,t$. Resulta
\begin{equation}
  \boxed{\;|h_k(t)| \;\propto\; \exp\!\left[-\frac{\big(k\,|\omega'|\,\sigma\,t\big)^2}{4}\right]\;}
  \label{eq:gaussiana}
\end{equation}

Esto es \textbf{una gaussiana en el tiempo}, no una exponencial. En escala
logarítmica, $\log|h_k|$ es una \emph{parábola} en $t$, o equivalentemente una
\emph{recta} frente a $t^2$. Es la firma inconfundible del phase mixing, y la
usaremos como test en la sección~\ref{sec:autogravedad}.

Obsérvese además la dependencia en $k$: los armónicos altos decaen como
$\exp(-k^2\cdot\text{const}\cdot t^2)$, muchísimo más deprisa. La estructura fina
angular se borra antes que la gruesa.

\subsection{El ancho de franja, y su consecuencia numérica}

Una forma complementaria de ver lo mismo: en el plano $(\Qr,\Jr)$ una línea
inicialmente vertical se inclina, porque el desplazamiento angular $\omega(\Jr)t$
depende de $\Jr$. Al tiempo $t$, la separación en $\Jr$ entre dos franjas
consecutivas del armónico $k$ es
\begin{equation}
  \Delta \Jr = \frac{2\pi}{k\,|\omega'|\,t} .
  \label{eq:franja}
\end{equation}

\begin{idea}
\textbf{La consecuencia práctica más importante de estas notas.} La estructura se
afina como $1/t$. Cualquier método numérico que represente $F$ con resolución
finita en $\Jr$ dejará de resolver las franjas a partir de cierto tiempo. Para
llegar a $t$ hace falta una resolución en acción
\[
  N_J \;\gtrsim\; \frac{4\,k\,\Delta\omega\,t}{2\pi},
\]
donde $\Delta\omega$ es el rango de frecuencias que abarca la distribución. El
requisito \textbf{crece linealmente con el tiempo final}, y por debajo de él el
resultado no es mediocre: es basura, por aliasing. Es el mismo fenómeno que la
\emph{recurrencia} conocida en simulaciones de plasmas~\cite{birdsall2004}.
\end{idea}

\section{Amortiguamiento de Landau: en qué se diferencia}
\label{sec:landau}

Todo lo anterior vale para partículas de prueba en un potencial fijo. Si el
sistema es autogravitante, la perturbación de la densidad genera una perturbación
del potencial, que a su vez actúa sobre las órbitas. Linealizando Vlasov--Poisson
alrededor de un equilibrio $F_0(\Jr)$ se obtiene una relación de dispersión cuyas
soluciones $\omega = \omega_r + \ii\gamma$ son los modos del sistema.

Landau~\cite{landau1946} demostró, en el contexto del plasma, que la respuesta
contiene una contribución de las partículas \emph{resonantes} —aquellas cuya
frecuencia orbital coincide con la del modo, $\omega(\Jr_{\rm res}) = \omega_r/k$—
que produce un decaimiento
\begin{equation}
  |h_k(t)| \;\propto\; \ee^{-\gamma t},
  \qquad
  \gamma \;\propto\; \left.\frac{\partial F_0}{\partial \Jr}\right|_{\Jr_{\rm res}} .
  \label{eq:landau}
\end{equation}
El signo de la derivada de la distribución en la resonancia decide si el modo se
amortigua o crece. En un sistema gravitatorio el tratamiento análogo es el de los
modos de van Kampen~\cite{vankampen1955} y Case~\cite{case1959}, y su versión
estelar está desarrollada en~\cite[\S5.5]{binney2008}.

\begin{idea}
\textbf{El test operativo.} Los dos mecanismos dejan firmas distintas y medibles:
\begin{center}
\begin{tabular}{lll}
\toprule
 & envolvente & en escala log \\
\midrule
phase mixing      & $\exp[-(k|\omega'|\sigma t)^2/4]$ & recta frente a $t^2$ \\
Landau damping    & $\exp(-\gamma t)$                 & recta frente a $t$ \\
\bottomrule
\end{tabular}
\end{center}
Se ajustan las dos y se comparan los coeficientes de determinación. No hace falta
nada más sofisticado.
\end{idea}

Hay una diferencia conceptual que conviene subrayar. El phase mixing es
cinemática: ocurre aunque las partículas no se vean entre sí. El amortiguamiento
de Landau requiere el acoplamiento colectivo y \emph{desaparece} si se apaga la
autogravedad. Por eso, en la práctica, la forma de aislarlo es correr el mismo
problema con y sin autogravedad y comparar. Ese control aparecerá una y otra vez.
\section{El método numérico}
\label{sec:metodo}

\subsection{Particle-in-cell}

La ecuación de Vlasov es una PDE en seis dimensiones (aquí, tras la reducción, en
dos). Discretizarla en malla es caro y difunde. La alternativa estándar es el
método \emph{particle-in-cell}: se representa $F$ por un conjunto de
\emph{macropartículas} que se mueven según las ecuaciones características, y el
campo se calcula depositándolas en una malla. La referencia canónica en plasmas es
Birdsall y Langdon~\cite{birdsall2004}, y en gravitación Hockney y
Eastwood~\cite{hockney1988}.

El ciclo de cada paso es:
\begin{enumerate}
  \item depositar las partículas en la malla radial para obtener $\rho(r)$;
  \item resolver Poisson y obtener $\Phi(r)$ y la fuerza;
  \item interpolar la fuerza de vuelta a cada partícula;
  \item avanzar las partículas un paso de tiempo.
\end{enumerate}

El depósito y la interpolación usan funciones de forma $S$ construidas con
B-splines de orden $n$, lo que suaviza el ruido de discretización a costa de
resolución. Es importante usar la \emph{misma} función en el depósito y en la
interpolación: esa simetría es lo que evita la autofuerza, es decir, que una
partícula sienta su propio campo.

\subsection{Integradores simplécticos}

El sistema es hamiltoniano, y eso debe respetarse. Un integrador genérico
(Runge--Kutta, por ejemplo) introduce una deriva secular en la energía que a
tiempos largos contamina cualquier medida. Un integrador \emph{simpléctico}
preserva exactamente la forma simpléctica, y como consecuencia conserva un
hamiltoniano modificado $\tilde H = H + \mathcal{O}(\Delta t^p)$; el error de
energía queda \emph{acotado} en vez de crecer~\cite{hairer2006}.

El más simple es el salto de rana en su forma \emph{kick--drift--kick}:
\begin{equation}
  \rp^{n+1/2} = \rp^{n} + \tfrac{\Delta t}{2}\,F(r^n),\quad
  r^{n+1}    = r^{n} + \Delta t\,\rp^{n+1/2},\quad
  \rp^{n+1}  = \rp^{n+1/2} + \tfrac{\Delta t}{2}\,F(r^{n+1}),
  \label{eq:leapfrog}
\end{equation}
que es de segundo orden. Para subir de orden, Yoshida~\cite{yoshida1990} demostró
que basta \emph{componer} pasos de salto de rana con coeficientes bien elegidos:
\begin{equation}
  \mathcal{S}_{p+2}(\Delta t) = \mathcal{S}_p(w_1\Delta t)\circ
                                \mathcal{S}_p(w_0\Delta t)\circ
                                \mathcal{S}_p(w_1\Delta t),
\end{equation}
con
\begin{equation}
  w_1 = \frac{1}{2-2^{1/3}},\qquad w_0 = -\frac{2^{1/3}}{2-2^{1/3}} .
  \label{eq:yoshida}
\end{equation}

\begin{aviso}
Algunos $w_i$ son \textbf{negativos}: esos subpasos van hacia atrás en el tiempo, y
eso es precisamente lo que cancela los términos de error de orden inferior. La
composición sigue siendo simpléctica en cada orden.

\medskip
\textsc{coste.} Subir de orden no es automáticamente ganancia. Para reducir el
error un factor $R$, un método de orden $p$ necesita $\Delta t/R^{1/p}$ y cuesta
$\sim n_{\rm etapas}\,R^{1/p}$. El de cuarto orden (3 etapas) es más barato que el
de sexto (7 etapas) hasta errores muy pequeños.
\end{aviso}

En este trabajo se usa el de cuarto orden con $\Delta t = 0.1$, es decir unos
$10^3$ pasos por periodo orbital.

\subsection{El diagnóstico $h_k$: cómo se evalúa}

La suma sobre partículas que aproxima~\eqref{eq:hk} es
\begin{equation}
  h_k \simeq 8\pi^2 L_0 \,\Delta r\,\Delta \rp
       \sum_{j=1}^{N} f_j\; B(\Jr_j)\;a_k\;\ee^{-\ii k \Qr_j},
  \label{eq:hksuma}
\end{equation}
donde $(\Qr_j,\Jr_j)$ se obtienen de $(r_j,p_{r,j})$ con el mapa de la
sección~\ref{sec:mapa} y $f_j$ es el peso de la partícula.

\begin{aviso}
\textsc{detalle numérico: la factorización.} El coeficiente angular de la función
de prueba,
\[
  a_k = \frac{1}{2\pi}\int_0^{2\pi} A(\Qr)\cos(k\Qr)\,\dd \Qr,
\]
no depende de las partículas. Calcularlo una vez fuera del bucle en lugar de una
vez por partícula redujo el coste total de la simulación en un factor 10. La
cuadratura se hace con Simpson sobre $[0,\pi]$ aprovechando la paridad; con 512
subintervalos el error es $2\times10^{-15}$ incluso para perfiles estrechos.
\end{aviso}

\subsection{Dos formas de colocar las partículas}
\label{sec:muestreo}

Este punto es, numéricamente, el más importante de todo el documento.

\paragraph{Monte Carlo.} Se sortean partículas con densidad proporcional a $F_0$
(por rechazo) y todas llevan el mismo peso. El error de cualquier promedio escala
como $N^{-1/2}$, sin umbral y sin meseta. Es robusto y no exige pensar.

\paragraph{Cuadratura.} Se colocan los nodos en una rejilla regular directamente en
$(\Qr,\Jr)$ y se les asigna como peso el valor de $F_0$. Como los nodos están
ordenados y los pesos son conocidos, esto no es muestreo estadístico sino una
\emph{regla de cuadratura}, y su error no tiene por qué obedecer a $N^{-1/2}$.

\begin{idea}
La diferencia es enorme y merece entenderse. En $\Qr$ el integrando es periódico y
suave, y la regla del trapecio periódica converge entonces \emph{más rápido que
cualquier potencia} de $1/N$~\cite{trefethen2014}: unos 25 nodos bastan. En $\Jr$
el integrando oscila como $\ee^{-\ii k\omega(\Jr)t}$, y ahí la convergencia tiene
\emph{umbral}: mientras no se cruce la condición de Nyquist de~\eqref{eq:franja} el
resultado es inútil; en cuanto se cruza, el error se desploma varios órdenes de
magnitud de golpe.

\medskip
La consecuencia práctica: el reparto de los $N$ nodos entre $\Qr$ y $\Jr$ importa
mucho más que su número total.
\end{idea}

\subsection{Poisson en simetría esférica}

Con simetría esférica la ecuación~\eqref{eq:poisson} se integra directamente. De
$\nabla^2\Phi = 4\pi G\rho$ se sigue, para la masa encerrada
$M(r)=4\pi\int_0^r r'^2\rho\,\dd r'$,
\begin{equation}
  \frac{\dd\Phi}{\dd r} = \frac{M(r)}{r^2}.
\end{equation}
El código integra el sistema con Runge--Kutta imponiendo $\Phi(0)=0$ y corrigiendo
después la constante para que se cumpla la condición asintótica $\Phi \to -M_{\rm
tot}/r$, es decir $\Phi + r\,\dd\Phi/\dd r \to 0$ en el borde.

\begin{aviso}
\textsc{cómo se verifica un solver de Poisson.} Fuera de la distribución la
solución debe ser exactamente kepleriana. Basta comparar la fuerza calculada con
$-M(r)/r^2$, con $M$ obtenida integrando la \emph{misma} densidad depositada. En
nuestro caso la razón entre ambas resultó $0.9927$, constante en todo el exterior;
el $0.7\%$ es el error de integrar por trapecios un perfil con un borde abrupto,
no del solver. Este tipo de comprobación cuesta minutos y descarta de un golpe toda
una clase de errores.
\end{aviso}
\section{Verificación del código}
\label{sec:verificacion}

Antes de usar un código para descubrir algo hay que demostrar que reproduce lo que
ya se sabe. La fórmula~\eqref{eq:hkexacto} permite calcular $h_k(t)$ con precisión
de máquina mediante cuadratura, y eso da un patrón contra el que medir.

\subsection{Las dos convergencias}

Con fondo fijo y $t_{\max}=2000$, el error medio de $h_k$ respecto del cálculo
exacto, normalizado por $h_0$:

\begin{center}
\begin{tabular}{lrrr}
\toprule
esquema & $N=500$ & $N=5000$ & $N=50000$ \\
\midrule
Monte Carlo ($k=1$)  & $4.5\times10^{-2}$ & $4.9\times10^{-3}$ & $2.3\times10^{-3}$ \\
Cuadratura ($k=1$)   & $2.7\times10^{-3}$ & $1.2\times10^{-9}$ & $1.3\times10^{-9}$ \\
\bottomrule
\end{tabular}
\end{center}

Las dos filas cuentan historias distintas. El Monte Carlo baja como $N^{-1/2}$ —el
ajuste sobre cinco realizaciones independientes por punto da pendientes entre
$-0.47$ y $-0.55$— de forma monótona y predecible. La cuadratura, en cambio,
\textbf{cae seis órdenes de magnitud entre dos puntos consecutivos}: eso es el
umbral de Nyquist de~\eqref{eq:franja} siendo cruzado. Con $N=500$ solo hay 20
nodos en $\Jr$, por debajo del requisito, y el resultado es peor que el del Monte
Carlo; con $N=5000$ hay 200, por encima, y el error cae al nivel del integrador.

\begin{aviso}
\textsc{una lección de método.} Con una sola realización por valor de $N$, las
pendientes del Monte Carlo salían entre $-0.29$ y $-0.65$: ruido de estimación
disfrazado de física. Hicieron falta cinco semillas por punto para que se
agruparan en torno a $-1/2$. Esto exigió, de paso, añadir un parámetro de semilla
al código: por defecto \texttt{gfortran} siembra su generador desde el sistema
operativo, de modo que dos corridas idénticas daban resultados distintos y
ninguna era reproducible.
\end{aviso}

\subsection{Pruebas agnósticas: ¿depende el acuerdo de la distribución elegida?}

Un peligro real al validar un código es que el acuerdo con la teoría refleje no
que el método es correcto, sino que el código se fue ajustando a un problema
concreto. La forma de descartarlo es cambiar el problema.

Se probaron tres distribuciones iniciales nuevas, cada una elegida para hacer una
\emph{predicción falsable} distinta:

\begin{itemize}
  \item \textbf{Bimodal}: dos grupos de acciones que mezclan a ritmos distintos, con
    parte angular $1+0.8\cos \Qr+0.3\cos 2\Qr$. Al ser un polinomio trigonométrico
    de grado dos, \textbf{no tiene armónicos por encima de $k=2$}: $h_3$ y $h_4$
    deben ser cero exactamente. Medidos: $2\times10^{-11}$ de $h_0$, ocho órdenes
    por debajo de la señal.
  \item \textbf{Espiral}: no separable, centrada en $\Qr=\beta \Jr$. Nace enrollada
    y se \emph{desenrolla} en $t^\ast=-\beta/\omega'(\Jr_0)$, donde $|h_k|$ pasa por
    un máximo. Predicho $t^\ast=720.6$; medido $722$, con la cadencia de salida de
    2 unidades.
  \item \textbf{King}: isoterma bajada, no suave —termina en una esquina. Su
    cuadratura converge como $h^2$ en vez de más rápido que cualquier potencia,
    exactamente como predice la teoría para un integrando solo continuo. Medido:
    factores 133 y 95 al multiplicar por diez los nodos, frente a los 100 de $h^2$.
\end{itemize}

Que un modo prohibido por la forma de $F_0$ salga a nivel de error numérico, y que
un pico de recoherencia aparezca en el instante calculado de antemano, son dos
pruebas que un código ajustado al problema anterior no pasaría por casualidad.

\section{Caso de estudio: qué pasa al encender la autogravedad}
\label{sec:autogravedad}

Aquí empieza la parte interesante, y es también donde se ve cómo se descarta una
explicación falsa. El montaje: la misma distribución gaussiana, con masa total
$a_0=10^{-3}$ —\emph{una milésima} de la del fondo isócrono—, ahora con su propia
gravedad encendida. $N=10^4$, cuadratura y Monte Carlo, $\Delta t=0.1$. Y, como
control imprescindible, la misma corrida con la autogravedad apagada.

\subsection{Resultado 1: la energía deja de conservarse a precisión de máquina}

\begin{center}
\begin{tabular}{lrr}
\toprule
 & deriva de $E$ & deriva de $h_0$ \\
\midrule
sin autogravedad          & $7.0\times10^{-11}$ & $-2.4\times10^{-10}$ \\
con autogravedad (cuad.)  & $5.0\times10^{-4}$  & $-1.602\times10^{-2}$ \\
con autogravedad (MC)     & $5.0\times10^{-4}$  & $-1.602\times10^{-2}$ \\
\bottomrule
\end{tabular}
\end{center}

La energía cae de $10^{-11}$ a $10^{-4}$. La razón es conceptual: al calcular la
fuerza depositando en una malla, \textbf{el esquema deja de ser exactamente
simpléctico}, porque la fuerza ya no deriva de un hamiltoniano de las posiciones
exactas. Ese $5\times10^{-4}$ es el suelo de confianza de estas corridas y hay que
tenerlo presente.

\subsection{Resultado 2: $h_0$ deriva, y eso es física}

$h_0$ deriva un $1.6\%$, treinta veces más que el error de energía, y
\emph{idéntico a tres cifras en los dos esquemas de muestreo}. No es ruido: es
física. Recuérdese que $h_0$ mide $\langle B(\Jr)\rangle$ y sólo es constante
mientras $\Jr$ lo sea. Al añadir el potencial propio, $\delta\Phi\sim a_0/r \sim
2\times10^{-4}$ frente a $|\Phi_{\rm iso}|\sim0.086$, las energías orbitales se
desplazan y con ellas las acciones.

\begin{idea}
\textbf{La deriva de $h_0$ es la medida directa de cuánto deja de conservarse
$\Jr$.} Es el diagnóstico más informativo de toda la simulación y es gratis.
Escala linealmente con la masa: $1.6\times10^{-3}$, $1.6\times10^{-2}$ y
$1.4\times10^{-1}$ para $a_0=10^{-4},10^{-3},10^{-2}$.
\end{idea}

\subsection{Resultado 3: sigue siendo phase mixing, no Landau}

Aplicando el test de la sección~\ref{sec:landau} al tramo inicial de decaimiento:

\begin{center}
\begin{tabular}{lrrl}
\toprule
 & $R^2$ frente a $t^2$ & $R^2$ frente a $t$ & veredicto \\
\midrule
sin autogravedad         & $1.00000$ & $0.938$ & gaussiana \\
con autogravedad (cuad.) & $0.99861$ & $0.938$ & gaussiana \\
con autogravedad (MC)    & $0.99859$ & $0.938$ & gaussiana \\
\bottomrule
\end{tabular}
\end{center}

La envolvente es gaussiana, no exponencial. Con una masa perturbadora de $10^{-3}$
el efecto colectivo no cambia la \emph{naturaleza} del decaimiento. Lo que sí hace
es \textbf{acelerarlo un 2.5\%}: el coeficiente del ajuste pasa de
$-3.726\times10^{-6}$ a $-3.820\times10^{-6}$ (cuadratura) y
$-3.843\times10^{-6}$ (Monte Carlo). Esa aceleración coherente entre dos esquemas
independientes es lo que queda del amortiguamiento colectivo en este régimen.

\subsection{Resultado 4: una meseta inesperada}

Y aquí aparece lo que motivó todo lo demás. Con autogravedad, $|h_1|$ deja de
decaer, toca un mínimo en torno a $t\simeq1200$, y se estabiliza en
$1.77\times10^{-3}$ de $h_0$. Sin autogravedad sigue bajando monótonamente.

La tentación es leerlo como un modo colectivo que sobrevive al phase mixing —sería
el amortiguamiento de Landau que uno busca—. Antes de creerlo hay que descartar lo
aburrido. Las dos hipótesis rivales predicen cosas opuestas:

\begin{center}
\begin{tabular}{lll}
\toprule
 & ruido de discretización & respuesta física \\
\midrule
al refinar la malla & cambia & no cambia \\
al subir $N$        & baja   & no cambia \\
al subir el orden del spline & baja & no cambia \\
al subir $a_0$      & no necesariamente & escala con $a_0$ \\
\bottomrule
\end{tabular}
\end{center}

Los cuatro barridos, con el valor de la meseta en unidades de $h_0$:

\begin{center}
\begin{tabular}{llllll}
\toprule
\textbf{malla} $\Delta r$ & $0.2$ & $0.1$ & $0.05$ & $0.025$ & \\
(con $\Delta t$ fijo)     & $1.771$ & $1.770$ & $1.770$ & $1.771$ & $\times10^{-3}$\\
\midrule
\textbf{orden B-spline}   & 1 & 2 & 3 & & \\
                          & $1.770$ & $1.770$ & $1.770$ & & $\times10^{-3}$\\
\midrule
\textbf{nodos en} $\Qr$   & 25 & 100 & 200 & & \\
                          & $1.7698$ & $1.7698$ & $1.7680$ & & $\times10^{-3}$\\
\midrule
\textbf{partículas} $N$   & $10^3$ & $10^4$ & $10^5$ & & \\
                          & $1.763$ & $1.770$ & $1.770$ & & $\times10^{-3}$\\
\midrule
\textbf{masa} $a_0$       & $10^{-4}$ & $10^{-3}$ & $10^{-2}$ & & \\
                          & $0.184$ & $1.771$ & $15.61$ & & $\times10^{-3}$\\
\bottomrule
\end{tabular}
\end{center}

\textbf{Nada de la discretización lo mueve} —ni ocho veces más resolución de malla,
ni B-splines cúbicos, ni ocho veces más nodos angulares, ni \emph{cien veces} más
partículas: tres dígitos iguales—. Y escala con la masa: la pendiente log-log vale
$+0.965$, frente al $+1$ de una respuesta lineal.

\begin{aviso}
\textsc{un error que merece quedar escrito.} En el primer análisis se concluyó lo
contrario: que la meseta era ruido amplificado. El argumento se apoyaba en el
escalado con $N$ \emph{del Monte Carlo}, que daba pendiente $-0.28$. Lo que fallaba
es que el nivel físico ($1.77\times10^{-3}h_0 = 4.4\times10^{-9}$) queda por debajo
del ruido del Monte Carlo en todo el rango explorado: sus mesetas valían
$5.9$, $2.4$ y $1.6\times10^{-8}$. El Monte Carlo estaba midiendo su propio ruido, y
ajustarlo a «ruido más una constante» era extrapolar desde tres puntos que aún no
habían tocado suelo. La cuadratura resuelve la señal ya con $N=10^3$. Correr los dos
esquemas no era redundancia: era lo que permitía distinguirlos.
\end{aviso}

\subsection{Pero no es un modo}

Establecido que es físico y lineal en $a_0$, la pregunta siguiente es qué modo es.
Un modo colectivo oscila: $h_k \sim \ee^{-\ii\omega_r t}$. Basta mirar la fase.

\begin{center}
\begin{tabular}{lrl}
\toprule
ventana & $\omega$ medida & \\
\midrule
temprana, $t\in[100,400]$   & $0.0607$ & coincide con $\omega$ orbital ($0.0629$) \\
tardía, $t\in[1300,2000]$   & $0.0008$ & \\
$t\in[15000,20000]$         & $5\times10^{-7}$ & \\
\bottomrule
\end{tabular}
\end{center}

En la ventana temprana la fase avanza a la frecuencia orbital, como debe ser: es el
continuo, $h_1\sim\ee^{-\ii\omega(\Jr_0)t}$. En la tardía \textbf{la fase está
congelada}. A lo largo de 200 periodos orbitales la amplitud pasa de
$1.7692\times10^{-3}$ a $1.7713\times10^{-3}$ y la fase se mantiene en $0.000$. Un
modo dentro de la banda~\eqref{eq:banda} habría acumulado más de $10^3$ radianes.

No es un modo. Es un residuo \emph{estático}.

\subsection{La explicación: son las coordenadas}

Un residuo estático y proporcional a $a_0$ sugiere una explicación que no es ni
ruido ni dinámica colectiva:

\begin{idea}
La simulación evoluciona en el potencial \emph{total} $\Phi_{\rm iso}+\delta\Phi$,
pero el diagnóstico calcula $(\Qr,\Jr)$ con el mapa del isócrono \emph{solo}. Una
distribución que ha mezclado por completo en las variables verdaderas —y que por
tanto depende únicamente de $\Jr_{\rm verdadero}$— \textbf{no} sale independiente
del ángulo al proyectarla con el mapa equivocado. El desajuste es
$\mathcal{O}(a_0)$, estático, y produce exactamente un $h_k$ residual constante.
\end{idea}

Esta hipótesis hace una predicción comprobable: si es cierta, la acción
\emph{isócrona} de una partícula individual no debe ser constante a lo largo de su
órbita real, sino \textbf{oscilar con el periodo orbital}. Guardando instantáneas
cada 4 unidades de tiempo y siguiendo partículas concretas:

\begin{center}
\begin{tabular}{rrrr}
\toprule
$\langle \Jr\rangle$ & oscilación pico a pico & $\omega$ de la oscilación & $\omega$ orbital \\
\midrule
$0.3000$ & $0.88\%$ & $0.0502$ & $0.0500$ \\
$0.4482$ & $0.79\%$ & $0.0439$ & $0.0426$ \\
$0.5914$ & $0.82\%$ & $0.0376$ & $0.0368$ \\
\bottomrule
\end{tabular}
\end{center}

\textbf{La acción isócrona oscila exactamente a la frecuencia orbital}, con
amplitud del orden del uno por ciento y una deriva neta diez veces menor. Es
decir: no hay difusión física de la acción; hay una coordenada que se bambolea una
vez por órbita porque no es la coordenada adecuada al potencial real.

Eso cierra el caso. La meseta no mide dinámica colectiva: mide el error sistemático
de usar el mapa ángulo--acción del potencial equivocado. Es física en el sentido de
que cuantifica una diferencia real entre dos potenciales, pero no dice nada sobre
modos ni sobre amortiguamiento.
\section{El régimen fuerte: construir el mapa ángulo--acción numéricamente}
\label{sec:mapanumerico}

La conclusión de la sección anterior deja una tarea clara: si la autogravedad no es
despreciable, hay que abandonar el mapa analítico del isócrono y construir el del
potencial real. Afortunadamente, con un grado de libertad esto es directo.

\subsection{Las cuadraturas}

Dado un potencial efectivo $\Phi_{\rm ef}(r)$ cualquiera y una energía $E$:

\begin{enumerate}
  \item \textbf{Puntos de retorno.} Resolver $E=\Phi_{\rm ef}(r)$ para $r_-<r_+$.
    Como $\Phi_{\rm ef}$ tiene un único mínimo, basta bisección más Newton.
  \item \textbf{Acción y periodo.}
    \begin{equation}
      \Jr(E)=\frac{1}{\pi}\int_{r_-}^{r_+}\!\!\sqrt{2(E-\Phi_{\rm ef})}\,\dd r,
      \qquad
      T_r(E)=2\!\int_{r_-}^{r_+}\!\frac{\dd r}{\sqrt{2(E-\Phi_{\rm ef})}} .
    \end{equation}
  \item \textbf{Ángulo.} Con $t(r)$ el tiempo transcurrido desde el pericentro,
    \begin{equation}
      \Qr = \frac{2\pi}{T_r}\int_{r_-}^{r}\frac{\dd r'}{\sqrt{2(E-\Phi_{\rm ef})}}
      \quad (\rp>0),
      \qquad
      \Qr = 2\pi - (\text{lo anterior}) \quad (\rp<0).
    \end{equation}
\end{enumerate}

Obsérvese que el mapa \emph{directo} $(r,\rp)\to(\Qr,\Jr)$, que es el que necesita
el diagnóstico, no requiere invertir nada: son tres cuadraturas. La inversión por
Newton solo hace falta para el mapa contrario, al colocar nodos de cuadratura.

\subsection{La sustitución que lo hace viable}

Los integrandos divergen como $(E-\Phi_{\rm ef})^{-1/2}$ en los puntos de retorno.
Una cuadratura ingenua converge como $\sqrt{h}$ y es inservible. El remedio
estándar es el cambio de variable
\begin{equation}
  r = \frac{r_++r_-}{2} + \frac{r_+-r_-}{2}\,\sin\theta,
  \qquad \theta\in[-\tfrac{\pi}{2},\tfrac{\pi}{2}],
  \label{eq:sustitucion}
\end{equation}
que anula la singularidad: cerca de $r_\pm$ se tiene $E-\Phi_{\rm ef}\propto
(r_\pm - r)$, y el jacobiano $\dd r = \tfrac{r_+-r_-}{2}\cos\theta\,\dd\theta$
aporta justo el factor que hace falta. Tras el cambio el integrando es suave y
Gauss--Legendre o el trapecio convergen espectralmente. El procedimiento está
descrito en~\cite[apéndice de \S3.5]{binney2008}.

\begin{aviso}
\textsc{coste, y cómo domarlo.} Hacer estas cuadraturas por partícula y por paso de
salida cuesta $\sim10^9$ evaluaciones para $N=10^4$ y mil salidas. La solución es
que, con $L$ fijo, $\Jr(E)$, $T_r(E)$ y $r_\pm(E)$ son funciones de \emph{una sola}
variable: se tabulan una vez sobre $\sim10^3$ puntos y se interpolan con splines.
Solo queda por partícula la integral parcial del ángulo, que es corta. El coste
baja a algo comparable al diagnóstico actual.
\end{aviso}

\subsection{¿Respecto de qué potencial?}

Con autogravedad $\Phi$ evoluciona, así que $\Jr$ solo es un invariante
\emph{adiabático}. Hay que elegir marco:

\begin{itemize}
  \item \textbf{Potencial instantáneo} $\Phi(r,t)$: $\Jr(t)$ es la acción
    instantánea, y sus cambios miden evolución genuinamente no adiabática.
  \item \textbf{Potencial medio o relajado} $\Phi_0(r)$: marco fijo, $\Jr$
    conservado para el campo medio, y $\delta F$ se analiza como perturbación
    alrededor de un equilibrio. \textbf{Es lo estándar} para el análisis de modos, y
    es el marco en el que la teoría de Landau está formulada.
\end{itemize}

Para estudiar phase mixing y amortiguamiento conviene lo segundo: dejar relajar el
sistema, tomar $\Phi_0$ como el potencial medio del estado final, construir el mapa
numérico para $\Phi_0$, y analizar $\delta F$ en esas variables.

\subsection{Cómo validarlo}

Hay una prueba unitaria evidente y muy fuerte: \textbf{aplicar el mapa numérico al
isócrono puro y comprobar que reproduce el analítico}. Si $\Jr_{\rm num}(E)$
coincide con~\eqref{eq:Jiso} a precisión de máquina, y $\Qr_{\rm num}$ con el mapa
de~\eqref{eq:kepler}, la implementación es correcta.

Y hay una segunda prueba, específica de estas notas: \textbf{la meseta de la
sección~\ref{sec:autogravedad} debe desaparecer}. Si la interpretación es correcta,
al analizar con las variables del potencial total, $h_1$ debe decaer sin meseta. Es
una predicción tajante, y sería la confirmación definitiva.

\subsection{Dónde está el umbral, con números}

De lo medido: la deriva de $h_0$ —que \emph{es} la medida de cuánto deja de
conservarse $\Jr_{\rm iso}$— vale $1.6\%$ con $a_0=10^{-3}$ y $14\%$ con
$a_0=10^{-2}$. El mapa analítico aguanta bien hasta $\sim10^{-3}$ y empieza a
fallar de forma clara hacia $10^{-2}$. Ese es el punto a partir del cual el trabajo
de la sección~\ref{sec:mapanumerico} deja de ser opcional.

\section{Cuántas partículas hacen falta}
\label{sec:cuantas}

Una pregunta práctica recurrente, con respuesta distinta para cada esquema.

\paragraph{Monte Carlo, fondo fijo.} El error escala como $N^{-1/2}$ sin meseta.
De los datos, $\varepsilon \simeq 3\times10^{-2}\sqrt{500/N}$, luego
\begin{equation}
  N \simeq 500\left(\frac{3\times10^{-2}}{\varepsilon}\right)^{2}.
\end{equation}
Un 1\% cuesta $\sim5\times10^3$ partículas; un 0.1\%, $\sim5\times10^5$. Cada
dígito adicional cuesta cien veces más.

\paragraph{Cuadratura, fondo fijo.} Aquí $N$ es la variable equivocada: lo que manda
es el reparto. Se necesita
\begin{equation}
  N_J \gtrsim \frac{4\,k_{\max}\,\Delta\omega\,t_{\max}}{2\pi},
  \qquad
  N_Q \sim 25\text{--}100 \ \text{según lo estrecho que sea } F_0 \text{ en } \Qr .
\end{equation}
Para $t_{\max}=2000$, $k\le4$ y $\Delta\omega\simeq0.02$ eso da $N_J\gtrsim107$, y
con $N_Q=25$ el total es $N\gtrsim2.5\times10^3$. \textbf{El requisito crece lineal
con $t_{\max}$}: es el único parámetro que no perdona.

\paragraph{Caso autogravitante.} Es donde la diferencia se vuelve dramática. La
cuadratura resuelve la señal colectiva con $N=10^3$; el Monte Carlo seguía
dominado por su ruido con $N=10^5$. \textbf{Elegir bien el esquema valió más que
multiplicar por mil el número de partículas.}

\section{Lecciones de método}

Más allá de la física, el desarrollo dejó algunas lecciones que probablemente se
generalicen.

\begin{enumerate}
  \item \textbf{Un control es imprescindible.} Casi todas las conclusiones de la
    sección~\ref{sec:autogravedad} descansan en comparar con la misma corrida sin
    autogravedad. Sin ella, la deriva de $h_0$ y la meseta habrían sido
    inatribuibles.

  \item \textbf{Cuando dos hipótesis predicen signos opuestos, hay un experimento.}
    Refinar la malla \emph{baja} el error de fuerza y \emph{sube} el ruido de
    disparo. Un solo barrido decide entre las dos.

  \item \textbf{Mirar la fase, no solo el módulo.} La distinción entre «modo
    colectivo» y «residuo estático» se resolvió en una línea al desenrollar la fase.
    Un espectro de Fourier sobre un registro de ocho periodos no lo habría visto.

  \item \textbf{Dos discretizaciones independientes valen más que una muy fina.} El
    Monte Carlo y la cuadratura fallan de maneras distintas; cuando coinciden a tres
    cifras, el resultado es sólido, y cuando discrepan, la discrepancia misma
    informa.

  \item \textbf{Los resultados negativos hay que escribirlos.} La primera
    interpretación de la meseta era falsa. Dejar constancia del porqué evita que
    alguien —empezando por uno mismo— repita el razonamiento.
\end{enumerate}

\section{Para seguir}

Lo que queda pendiente, en orden de dificultad:

\begin{itemize}
  \item Implementar el mapa ángulo--acción numérico de la
    sección~\ref{sec:mapanumerico} y comprobar que la meseta desaparece.
  \item Con ese mapa, explorar $a_0\gtrsim10^{-2}$, donde el amortiguamiento
    colectivo deja de ser una corrección del $2.5\%$ y puede llegar a dominar.
  \item Buscar modos discretos genuinos: resolver el problema de autovalores
    linealizado, o medir con cuidado $\omega_r$ y $\gamma$ en corridas largas, y
    comprobar si caen dentro o fuera de la banda~\eqref{eq:banda}. Un modo dentro de
    la banda es resonante y por tanto amortiguado por Landau; fuera, no.
  \item Levantar la restricción de $L$ fijo, lo que reintroduce un segundo grado de
    libertad y con él resonancias entre las dos frecuencias.
\end{itemize}

\begin{thebibliography}{99}

\bibitem{binney2008}
J.~Binney y S.~Tremaine,
\emph{Galactic Dynamics}, 2\textsuperscript{a} ed.,
Princeton University Press, 2008.
\newblock El texto de referencia. El capítulo~3 desarrolla las variables
ángulo--acción y el potencial isócrono; el~5, la teoría de perturbaciones y el
amortiguamiento de Landau en sistemas estelares.

\bibitem{henon1959}
M.~Hénon,
\emph{L'amas isochrone I},
Annales d'Astrophysique \textbf{22} (1959) 126.
\newblock Introduce el potencial isócrono y demuestra que sus acciones son
expresables en forma cerrada.

\bibitem{landau1946}
L.~D.~Landau,
\emph{On the vibrations of the electronic plasma},
J. Phys. USSR \textbf{10} (1946) 25.
\newblock El artículo original. La versión en inglés está recogida en
\emph{Collected Papers of L. D. Landau}, Pergamon, 1965.

\bibitem{vankampen1955}
N.~G.~van Kampen,
\emph{On the theory of stationary waves in plasmas},
Physica \textbf{21} (1955) 949.

\bibitem{case1959}
K.~M.~Case,
\emph{Plasma oscillations},
Annals of Physics \textbf{7} (1959) 349.
\newblock Junto con el anterior, la interpretación del amortiguamiento de Landau
como superposición de modos del continuo que se desfasan, es decir, como phase
mixing en el espacio de velocidades.

\bibitem{lyndenbell1967}
D.~Lynden-Bell,
\emph{Statistical mechanics of violent relaxation in stellar systems},
Monthly Notices of the Royal Astronomical Society \textbf{136} (1967) 101.
\newblock El mecanismo de relajación colectiva rápida; complementa al phase mixing
para perturbaciones grandes.

\bibitem{birdsall2004}
C.~K.~Birdsall y A.~B.~Langdon,
\emph{Plasma Physics via Computer Simulation},
Taylor \& Francis, 2004 (edición original McGraw-Hill, 1985).
\newblock La referencia clásica de PIC. Discute funciones de forma, autofuerza y el
fenómeno de recurrencia.

\bibitem{hockney1988}
R.~W.~Hockney y J.~W.~Eastwood,
\emph{Computer Simulation Using Particles},
Adam Hilger, 1988.

\bibitem{yoshida1990}
H.~Yoshida,
\emph{Construction of higher order symplectic integrators},
Physics Letters A \textbf{150} (1990) 262.
\newblock Los coeficientes de composición usados aquí.

\bibitem{hairer2006}
E.~Hairer, C.~Lubich y G.~Wanner,
\emph{Geometric Numerical Integration: Structure-Preserving Algorithms for
Ordinary Differential Equations}, 2\textsuperscript{a} ed.,
Springer, 2006.
\newblock Por qué los integradores simplécticos acotan el error de energía en vez
de dejarlo crecer: análisis del hamiltoniano modificado.

\bibitem{trefethen2014}
L.~N.~Trefethen y J.~A.~C.~Weideman,
\emph{The exponentially convergent trapezoidal rule},
SIAM Review \textbf{56} (2014) 385.
\newblock La razón de que bastan $\sim25$ nodos en el ángulo: para integrandos
periódicos y analíticos el trapecio converge más rápido que cualquier potencia.

\end{thebibliography}

\end{document}
